\section{Discussion}
\label{sec:discussion}

The project kept a lab notebook (\texttt{docs/research/lab\_notebook.md}) in which every defect,
correction and prediction was written down when it happened and never edited afterwards; later
corrections are new entries. This section is built from it. Its value is not that the benchmark
is good, but that the ways a careful process still went wrong are recorded, and several recur.

\subsection{Recurring failure patterns}

\paragraph{A control measured on the inputs is not a control on the system.} The most frequent
pattern, recorded under several names (``a guard with two doors'', ``a fix that keeps the old
behaviour as a floor''). A parameter digest guarded the generator's inputs while the claim lived
in its output: the draw could change with every parameter identical, which is why datasets are now
identified by a fingerprint over their rows. A statistical band kept a fixed floor that, at release
scale, was several times wider than the band it was meant to bound, so the check could not fail.
Out-of-fold encoding removed each row's own label but left the other rows of the same fraud
incident in the estimate, inflating every single-feature AUC by about \nSiblingLeak{} until folds
were grouped by account; later, account grouping turned out not to protect features whose unit of
history is a counterparty or a map cell. The single-column separation ceiling was true of the columns it
measured and false of the engineered features it was quoted about. The fingerprint fix itself first hashed only some of
the columns. In each case the obvious form of the defect (the individual value, the row) was
fixed and the aggregate form (the set, the group, the regime) was not. The habit that follows,
now an exit criterion: for every guard, name the object it ranges over and read the sentence it
is supposed to justify; if the nouns differ, it does not justify the sentence. And prefer a guard
with no scope parameter to one with a well-chosen scope.

\paragraph{An asserted protection is worse than a missing one.} A documented cross-account control
that did not in fact control anything, and a code comment explaining why a truncated corpus could
not bias a feature, both survived review by being reassuring rather than invisible. The comment's
argument turned out to be right for the feature it named, and was only established once it was
measured (Section~\ref{sec:floor}); the search it prompted found three further features with the
flaw the comment had ruled out for its own. Controls must now name the mutation that would detect their failure, and the
rule recorded is that an aggregate over a window can be derived from a corpus, while a statement
about all of history cannot.

\paragraph{A test that cannot fail proves nothing.} Fixtures built to be small rather than to
contain the condition under test passed vacuously: a fold test whose categories were pure, a drift
test whose fixture had no drifting rows. Tests now assert their precondition before their
conclusion, and fixtures are specified by the property they must have before their size.

\paragraph{The wrong denominator.} The central methodological error of the dataset milestone was
standardising a detector's offset against a theoretical null rather than the observed seed-to-seed
spread; a
single three-sigma reading turned out to be noise once repeated seeds were run. The first model margin
reported subtracted a feature measured on a few hundred rows holding two frauds from a model
measured on the full test set. Every baseline now prints the rows and fraud it rests on, and a
partial-coverage feature is never subtracted from a model.

\paragraph{Records that run ahead of the facts.} A status table marked a gate as met by citing a
run that did not exist; an evidence file attributed to a commit had been produced from a dirty
working tree; a milestone tag was placed on a commit whose continuous integration had not passed.
Each is a hand-written summary of a machine-checkable fact, drifting in the flattering direction.
Status is now derived from the existence of evidence artefacts, evidence files carry the tree
state they were produced from, and a tag requires a green run on the exact commit.

\paragraph{Knowing a failure mode does not immunise against it.} Several of these patterns recurred
after they had been analysed and recorded. The lesson
the notebook draws is that a defence that depends on remembering a rule at the moment of temptation
is not a defence; each lesson is useful only with the question ``what would enforce this?''
attached, and most of the rules above are now tests.

\paragraph{A hedge is often an unmeasured claim.} Several caveats (``the delay channel might be
mostly an artefact'', ``burstiness is not country-specific'') turned out to be measurable, and
measuring them replaced a hedge with a result: about \nEventDelayShare{} of the delay channel was
artefact, and leave-one-country-out is null. The test for keeping a caveat is to ask what
measurement would remove it.

\subsection{Pre-registered predictions}

The notebook recorded \nPredTotal{} falsifiable predictions before the evidence existed, and
reports each outcome where the prediction was made, whether confirmed or refuted
(Table~\ref{tab:predictions}).

\begin{table}[htbp]
\centering
\small
\caption{Pre-registered predictions and outcomes, from the lab notebook.}
\label{tab:predictions}
\begin{tabular}{p{0.52\linewidth}p{0.38\linewidth}}
\toprule
Prediction & Outcome \\
\midrule
Shortcut-detector power dips when the eighth scenario appears above the minimum scale &
Refuted: power rose monotonically, no dip \\
The detector's small positive offset is real and constant with scale &
Refuted: pooled over seeds, the interval straddles chance \\
A counterparty-keyed feature will force an eighth field into the feature contract &
Confirmed on mechanism, wrong on sequence \\
\texttt{accounts\_per\_device\_7d} cannot be made fold-safe without an owner decision &
Refuted: grouping accounts that share a device resolves it \\
A long-tailed takeover lead lowers the event-delay channel to
\nEventDelayBandLow{}--\nEventDelayBandHigh{} &
Confirmed: \nEventDelayAfter{}, at the top of the band \\
The reversal scam is caught at \nRevPredLow{}--\nRevPredHigh{} recall &
Refuted: \nRevRecall{} \nRevCI{}; below the range for the point estimate, not for the whole
interval \\
\bottomrule
\end{tabular}
\end{table}

Of the first \nPredEarly{}, \nPredEarlyRight{} was partly right. The claim this supports is not that
the author predicted well. It is that falsifiable predictions were recorded before the evidence,
that most were wrong, that each refutation cost a paragraph rather than a shipped result, and that
several of the guards now in the system exist because of what the refutations found.

\subsection{A benchmark audit protocol}

The checks below are what we would run before trusting any synthetic benchmark, whoever built it.
Each one exists because its absence hid something here: Table~\ref{tab:protocol} states what each
check detects and what it found on FraudShield-EAC. They are cheap: every one is a measurement on
data that already exists, and together they decide which claims the benchmark can carry.

\begin{table}[htbp]
\centering
\small
\caption{A benchmark audit protocol: the checks, what each detects, and what each found on
FraudShield-EAC. Every finding is reported in the section named.}
\label{tab:protocol}
\begin{tabular}{p{0.26\linewidth}p{0.3\linewidth}p{0.36\linewidth}}
\toprule
Check & What it detects & What it found here \\
\midrule
The best single \emph{engineered} feature, not only the best raw column &
A benchmark one rule can solve, and model metrics quoted against chance &
The velocity ratio alone reaches \nFloorM{}, while the strongest raw column reaches
\nColumnMaxAUCB{}; every metric is now a margin over the floor (Section~\ref{sec:floor}) \\
\addlinespace
Leave-one-group-out \emph{and} one-group-only &
Redundancy read as irrelevance: the first check alone cannot tell them apart &
Removing any group costs at most \nGroupMaxDrop{}, yet four groups each reach at least
\nGroupGeographic{} alone (Section~\ref{sec:redundancy}) \\
\addlinespace
Does each generalisation experiment withhold a \emph{mechanism}, or only a label? &
Experiments that cannot fail, reported as evidence of generalisation &
No fraud parameter is keyed by country, so leave-one-country-out costs at most \nLocoMaxLoss{};
the variant held out in time is still a burst and is caught at \nNovelRecall{}
(Section~\ref{sec:generalisation}) \\
\addlinespace
A pre-registered variant from an independently documented typology &
Whether detection survives without the benchmark's dominant structure &
Recall \nRevRecall{} \nRevCI{} against \nBaseVarRecall{} \nBaseVarCI{}, below the registered
prediction (Section~\ref{sec:generalisation}) \\
\addlinespace
A fingerprint over the output rows, beside any digest of the parameters &
Reports that silently describe a different draw &
A refactor that changed no parameter value re-drew the whole benchmark; datasets are now
identified by their rows (Section~\ref{sec:repro}) \\
\addlinespace
Encodings and folds grouped by each feature's unit of history &
Leakage between rows of one incident, which inflates every feature's score &
Out-of-fold encoding grouped by row inflated single-feature AUC by \nSiblingLeak{}; grouping by
account fixed it, and features keyed by counterparty or map cell needed their own grouping \\
\addlinespace
For every control, a named mutation that would kill it &
Controls and tests that cannot fail &
Vacuous fixtures (a fold test with pure categories, a drift test with no drift) and a documented
cross-account control that controlled nothing \\
\bottomrule
\end{tabular}
\end{table}

On FraudShield-EAC each of these changed what could be claimed.
